% Transcription — fonds Alexandre Grothendieck, shelfmark 64, batch 1 (pages 1-20).
% Established from the facsimile archives/batches/64/batch-01.pdf (a copy of 64.pdf
% fetched through the project relay, no cover sheet: PDF page k = archivists' page k),
% per the skill transcribe-grothendieck. The folder is 156 pages in eight batches;
% this is the first.
%
% Pass: Opus 5.5 (claude-opus-5-5), 2026-09-27 — first pass, unchecked against the pages by a human.
%
% Read in two parts in the same session: pages 2-13 by the main pass, page 1 and
% pages 14-20 by a parallel pass on the same model, then merged.
%
% Pages skipped: none. Page 1 is a cover in his hand (large roman « I », title
% « Fourbis combinatoires (ens. à 4 éléments, cubes, octaèdres, modules libres de
% rang 2 sur Z/4Z) »): no \page{1}, his words become the top heading; his numbered
% sections 1-3 are therefore set as subsections.
%
% His pagination: the even archive pages 4-20 carry his numbers 2-10
% (p. 4 = 2, p. 6 = 3, ..., p. 20 = 10); p. 2 is his first page, unnumbered.
%
% What the batch is about: « Fourbis combinatoires », the combinatorics behind
% level-2 structures. § 1: a 4-element set I is canonically an affine plane over
% F_2; V the translations, J = V^*, 1 -> V -> S_I -> S_J -> 1, Sylow subgroups,
% (Ens_4) -> (Gr) fully faithful, the elements of S_4 named after rotations of
% the cube, orientation sets omega(I) ~ omega(J). § 2: the combinatorial cube,
% in five equivalent descriptions ((S, A, R), diagonals D = S/a, codiagonals
% C = A/a, torsors under Gamma); its automorphism group D of order 8, the
% subgroups V_S ~ F_2^C, V_A ~ F_2^D, D^+ cyclic of order 4, the relation
% omega_Q ^ C ^ D ~ 1 and the conventions it depends on, the dual cube;
% Aut D, cubes as torsors under D_4, a boxed « Formulaire ». § 3 (from p. 19):
% the 2-subgroups of S_4, begun as a classification of all its subgroups.
%
% Notation fixed once for the batch: \mathfrak{S}_I symmetric groups;
% \mathfrak{D} = Aut Q (fraktur) against D the set of diagonals (roman);
% \mathbb{D}_4 = Aut(Q_0) for the pinned cube; W = \mathfrak{D}_{ab};
% \underline{\omega} orientation sets; \underline{a} the antipodism; V^* the
% nonzero elements of V (his *).
%
% \ill{}: 123. \uncertain{}: 55.
%
% Points to check against the page: p. 6, x -> Cent_G(x) lands on the cyclic
% subgroups of order 4 as written; p. 9, D^+_ab vs D_ab; p. 10, the last
% enumeration reads « B, S, omega_Q » where D, C are meant; p. 15, « D -> D^+ »
% with a clear +; p. 19, the bottom node of the triangle; p. 20 d), V^* where
% the count points to V. All transcribed as written, with notes.
%

\documentclass[11pt,a4paper]{article}
\input{../preamble/grothendieck}

\folder{64}
\batch{1}
\pages{1}{20}
\dating{[à partir de 1980- 1982]}
\watermark{Édition de démonstration}
\foldertitle{Modules courbes elliptiques en caractéristiques ≠2 : notes manuscrites (s.d.)}

\begin{document}

\section{I. Fourbis combinatoires (ens.\ à 4 éléments, cubes, octaèdres, modules libres de rang 2 sur $\mathbb{Z}/4\mathbb{Z}$).}

\note{titre pris sur la page~1, feuillet par ailleurs blanc qui sert de
couverture à la suite : il est de la main de l'auteur, en haut à droite,
sous un grand chiffre romain « I ». À gauche, un « (40) » entre parenthèses,
d'une écriture plus petite, sans doute un compte de pages ; il n'est pas sûr
qu'il soit de l'auteur.}

\subsection{1. Ens. à quatre éléments.}

\page{2}
Un ens. à quatre éléments $I$ est muni canoniquement d'une structure
d\struck{e torseur sous un \ill{}} \add{d'espace affine} de dimension 2
sur $\mathbb{F}_2$, i.e. d'un sous-groupe $V \subset \mathfrak{S}_I$,
isomorphe (non canon\supplied{iquement}!) à $\mathbb{F}_2^2$, tel que l'action
de $V$ sur $I$ en fasse un torseur. Le sous-groupe $V \subset \mathfrak{S}_I$
est le seul sous-groupe distingué abélien $\neq \{0\}$, \struck{\ill{}} il est
son propre \struck{normalisateur} \add{centralisateur} i.e.
$\mathfrak{S}_I/V$ opère fidèlement sur $V$, en fait on a
\[
  \mathfrak{S}_I/V \xrightarrow{\;\sim\;} \operatorname{Aut}_{\mathrm{gr}}(V)
  = \operatorname{Aut}_{\mathbb{F}_2\text{-vect}}(V)
  \xrightarrow{\;\sim\;} \mathfrak{S}_J
\]
i.e. $J = V^*$, la suite exacte
\[
  1 \to V \to \mathfrak{S}_I \to \mathfrak{S}_J \to 1
\]
\uncertain{étant} \struck{\uncertain{identique}} \add{\uncertain{identifiée}} à la suite
\[
  1 \to \underset{\text{translations}}{V} \to \operatorname{Aut}_{\mathrm{aff}}(I)
  \to \operatorname{Aut}_{\mathbb{F}_2}(V) \to 1 .
\]
\note{la parenthèse fermante après « $\operatorname{Aut}_{\mathrm{gr}}(V)$ »
et « $\operatorname{Aut}_{\mathbb{F}_2\text{-vect}}(V)$ » est écrite comme
un « $W)$ » ; la lecture $V$ est celle qu'impose le contexte. Le dernier terme
est surchargé : un mot biffé, encerclé, puis $\mathfrak{S}_J$.}

Les \add{trois} éléments \struck{\ill{}} de $J = V^*$ correspondent aux trois
partitions de type $(2,2)$ de $I$ (on peut \add{ainsi} dire que les élts de
$J$ correspondent aux droites \add{vectorielles} dans $V^*$, qui définissent
\struck{les} donc des \struck{r}relations d'équivalence dans $I$, vers les
espaces affines quotients de dim $1$) --- les \struck{\ill{}}classes
\uncertain{inversement} des \uncertain{l'ensemble} de ces \add{3} relations
d'équivalence correspondent aux 6 \struck{\uncertain{paires}}
\add{\uncertain{parties}} \uncertain{à} deux éléments de $I$) ou encore
\struck{aux} aux 6 droites affines dans l'espace affine $I$, l'application
\[
  \begin{aligned}
    \mathfrak{P}_2(I) &\longrightarrow \mathfrak{P}_{2,2}(I) \\
    A &\longmapsto \{A, I \setminus A\}
  \end{aligned}
\]
\marginal{\uncertain{ou aussi aux trois} translations \ill{} sans pts fixes}
\note{la note marginale, écrite de biais dans la marge gauche, se rattache à
« les éléments de $J$ correspondent aux trois partitions de type $(2,2)$ ».
La phrase du corps, depuis « --- les classes », est d'une lecture peu sûre ;
le sens mathématique en est clair (les six parties à deux éléments de $I$,
ou ses six droites affines, s'envoient deux à deux sur les trois partitions
de type $(2,2)$).}

\page{3}
s'identifiant à l'application qui associe à une droite affine de $I$ sa
\emph{direction}, qui est une droite de l'espace des translations $V$. La
fibre d'un \struck{\ill{}} $i \in \mathfrak{P}_{2,2}(I)$, i.e. l'un des
\struck{\ill{}} deux classes \add{de cette partition}, est dans l'espace
quotient $I/(D_i = \{0, \ldots\})$.\note{le second élément de $D_i$ est surchargé et n'est pas lu.}

Le foncteur « oubli de la structure » qui va de la catégorie des espaces
affines de dim $2$ sur $\mathbb{F}_2$, vers la cat. des ensembles à quatre
éléments
\[
  \mathrm{Aff}_2(\mathbb{F}_2) \xrightarrow{\;\approx\;} (\mathrm{Ens}_4)
\]
est une équivalence de catégories --- et \struck{\ill{}}
\add{i.e., si on définit la notion d'espace affine sur $\mathbb{F}_2$ par la
donnée d'une structure \struck{\ill{}} sur un $I$ \struck{tel} seul (donnée
d'un sous-groupe de $\mathfrak{S}_I$ par ex.)} Je dis que le foncteur
\[
  \begin{aligned}
    (\mathrm{Ens}_4) &\longrightarrow (\mathrm{Gr}) \\
    I &\longmapsto \mathfrak{S}_I
  \end{aligned}
\]
est pleinement fidèle, on en conclut que pour $\forall\, I$,
$\mathfrak{S}_I$ est un centre trivial et tout automorphisme de
$\mathfrak{S}_I$ est intérieur. \struck{Pourrais} \add{(\uncertain{il}}
suffirait bien sûr de l'énoncer pour $\mathfrak{S}_4$ ?). On va définir pour
ceci un foncteur \ill{}
\[
  (\mathrm{Gr})' \longrightarrow (\mathrm{Ens}_4)
\]
sur l'image essentielle $(\mathrm{Gr})'$ des groupes isomorphes à
$\mathfrak{S}_4$. Si \ldots\ est un tel groupe $G$, on considère l'ens.
\struck{\ill{}} \add{des} \struck{\ill{}} $3$-sous-groupes de Sylow, il y en a
exactement quatre ! Pour préciser ce point et \uncertain{récapituler}
\marginal{\ill{} fidèles \ill{}}
\note{le mot qui précède le foncteur $(\mathrm{Gr})' \to (\mathrm{Ens}_4)$
commence par un V majuscule (« V\ldots-à-\ldots ») et n'est pas lu ; le
sens est celui d'un foncteur en sens inverse du précédent. La note marginale
gauche, écrite de biais et soulignée, se rattache au foncteur
« oubli de la structure » ; seul « fidèles » s'y lit. Le « \ldots » après
« Si » est de l'auteur. Devant $(\mathrm{Gr})'$, un « $(\mathrm{Ens})$ » biffé.}

\page{4}\note{p.~2 de l'auteur, chiffrée en tête de page ; il ne numérote
que les pages paires de l'archive (ici 4, 6, \ldots, 20 portent ses
numéros 2 à 10), la page 2 étant sa première.}
la structure de torseur sur cet ens., on considère le sous-groupe abélien
distingué \struck{\ill{}} unique $V$ de $G$, on sait que $V$ est un
vectoriel de rang 2 sur $\mathbb{F}_2$ et
$G/V \xrightarrow{\;\sim\;} \operatorname{Aut}_{\mathbb{F}_2}(V)$, et on
sait que les \add{$3$-}sous-groupes de Sylow de $G$ correspondent aux
« relèvements » de l'unique $3$-sous-groupe de Sylow de
$G/V \simeq \mathfrak{S}_J$ ($J = V^*$), savoir
$\mathfrak{S}_J^+ = [\mathfrak{S}_J, \mathfrak{S}_J]$, d'indice 2 dans
$\mathfrak{S}_J$. \struck{\ill{}} L'ens. $I_G$ des $3$-sous-groupes de Sylow
de $G$ s'identifie donc à l'ens. des scindages de l'extension
\[
  1 \to V \to G^+ \to \mathfrak{S}_J^+ \to 1
\]
du $3$-groupe $\mathfrak{S}_J^+$ par le $2$-groupe $V$. Sur l'ens. de ses
scindages, $V$ opère donc transitivement par automorphismes intérieurs,
d'ailleurs $V^{\mathfrak{S}_J^+} = \{0\}$ ($\mathfrak{S}_J^+$ est le groupe des
permutations circulaires sur $V^* = J$), donc $I_G$ est bien un torseur sous
$V$.

Quand on part de $I$ et \struck{\ill{}} $G = \mathfrak{S}_I$, on trouve une
bijection
\[
  \begin{aligned}
    I &\longrightarrow I_G \\
    i &\longmapsto \mathfrak{S}^+_{I \setminus \{i\}}
      = 3\text{-ss-gr.\ de Sylow de } \mathfrak{S}_I
        \text{ qui fixe } i \,;
  \end{aligned}
\]
Ainsi \struck{les deux foncteurs} \add{\ill{}} foncteurs
$I \mapsto \mathfrak{S}_I$, $G \mapsto I_G$, et le composé
$I \to I(\mathfrak{S}_I)$, on a un isom\supplied{orphisme} \ill{} \struck{l'identité}
le foncteur identique,
\note{dans la dernière formule, « unique » est biffé à l'intérieur de la
formule ; la dernière ligne de la page, qui se poursuit page~5, n'est lue
qu'en partie.}

\page{5}
Pour le composé dans l'autre sens, on trouve que $G$ opère sur l'ens.\ $I_G$
de ses \ill{} ss-groupes de Sylow par automorphismes intérieurs, d'où
\[
  G \xrightarrow{\;\sim\;} \mathfrak{S}_{I_G}
\]
qui est un iso.

\textbf{Éléments de $\mathfrak{S}_I$.}

\begin{enumerate}
  \item[a)] $1$
  \item[b)] Les \add{trois} éléments \struck{\ill{}} \add{de} $V^*$, qui forment
    une orbite $J$ sous $\mathfrak{S}_I$, et correspondant aux élém.\ de
    $\mathfrak{P}_{2,2}(I)$,
    \marginal{« translations » (\ill{}) trois axes [sym.\ p.r.\ \ill{} des
    cubes]}
  \item[c)] Les \ill{} 8 éléments d'ordre 3, se répartissant en
    \struck{quatre} \add{deux} \struck{torseurs} \ill{} de 4 qui sont des
    torseurs (à droite, à gauche) sous $V$, savoir les deux élts $\pm 1$ de
    $\mathfrak{S}_J^+$. L'ens.\ de ces éléments s'identifie ainsi :
    \[
      I \times \underline{\omega}, \quad \text{où } \underline{\omega}
      = \omega(J)\ \bigl(\simeq \omega(I)\bigr)
    \]
    par l'application qui associe à tout $\rho \in \mathfrak{S}_I$ d'ordre 3
    son unique pt fixe dans $I$ (ou encore, le ss-groupe de Sylow qu'il
    engendre), \underline{et} l'orientation qu'il détermine \struck{\ill{}} ou
    encore son image dans $\mathfrak{S}_J^+$ :
    \marginal{« rotations » (ayant un pt fixe et un seul, lequel \ill{} $I$
    \ill{} $V$) [« rotations-sommets » des cubes]}
  \item[d)] Les \struck{\ill{}} \add{6} éléments \add{d'ordre 2} \struck{\ill{}}
    de $\mathfrak{S}_I$ qui s'envoient dans $\mathfrak{S}_J^-$, en
    correspondance 1-1 avec l'ens.\ des \underline{transpositions} dans $I$
    $\mathfrak{P}_2(I)$ --- on sait que les \ill{} qui se groupent
    \add{\uncertain{3}} \struck{\ill{}} \struck{paires}, par la transposition
    $\tau_A$, $\tau_{A'}$ \struck{\ill{}} ($A' = I \setminus A$) est l'unique
    \add{autre} élément d'ordre 2 tel que $\tau_A \tau_{A'} \in V^*$.
    \marginal{« réflexions » (ayant \struck{\ill{}} deux pts fixes \ill{})
    [« rotations-arêtes » des cubes]}
  \item[e)] Les 6 éléments d'ordre 4 de $\mathfrak{S}_I$, correspondant aux
    6 ordres circulaires sur $I$. \uncertain{Pour}
    \marginal{[« rotations-faces » des cubes]}
\end{enumerate}
\note{les quatre notes marginales de cette page, écrites de biais dans la
marge gauche, donnent pour chaque classe d'éléments de $\mathfrak{S}_I$ son
nom dans le groupe des rotations du cube ; elles sont d'une autre encre,
plus fine. Au point c), le mot souligné devant « 8 » n'est pas sûr. Au
point d), la phrase après « on sait que » est d'une lecture incertaine ; le
sens est que les six transpositions se groupent en trois paires $\tau_A$,
$\tau_{A'}$ avec $\tau_A\tau_{A'} \in V^*$.}

\page{6}\note{p.~3 de l'auteur.}
\struck{\ill{}} on a $u^2 \in V^*$, donc on a une application
\[
  \mathrm{Circ}(I) \longrightarrow V^*
\]
de degré 2 \uncertain{sur} (qui définit donc une structure de
\underline{cube}, \ill{} $V^*$ comme ens.\ des directions des faces,
$\mathrm{Circ}(I)$ comme l'ens.\ des six faces --- en fait ce cube est
canoniquement orienté ; cf.\ plus bas.
\marginal{$\circledast$ On peut dire aussi que les \ill{} éléments d'ordre 4
\ill{} se répartissent en trois paires $\{u, u^{-1}\}$ \ill{} sous-groupes
\ill{} de $\mathfrak{S}_I = G$ \ill{} d'ordre 4 \ill{} $V^*$. \ill{} --- $I$
\uncertain{est l'ens.\ des quatre} diagonales des cubes,}

L'isom.\ canonique
\[
  (**) \qquad \underline{\omega}(I) \simeq \underline{\omega}(J)
\]
On a
\[
  \mathfrak{S}_{I,\mathrm{ab}} \xrightarrow{\;\sim\;} \mathfrak{S}_{J,\mathrm{ab}}
  \simeq \{\pm 1\}
\]
les isom.\ étant donnés par la signature, d'où la possibilité de définir un
iso.\ \struck{canonique} \add{fonctoriel} $(**)$, \uncertain{moyennant}
cependant le choix d'une convention (deux choix possibles). Quand on fixe un
élément $i \in I$, alors on a
\[
  \underline{\omega}_i(I) \xrightarrow{\;\sim\;}
  \underbrace{\underline{\omega}(I \setminus \{i\})}_{\simeq\, \mathfrak{T}_{\mathrm{tot}}}
\]
\struck{\ill{}} $\ni$ l'ordre \add{total} \struck{\ill{}} de $I$ en associant à
un ordre \add{total} \uncertain{sur} $I \setminus \{i\}$ l'ordre donné, et qui
admet $i$ comme \underline{premier} élément (\uncertain{c'est là} la
convention \ill{} opposée aurait été de prendre $i$ comme \underline{dernier}
élément).
\note{sous l'accolade, « $\simeq \mathfrak{T}_{\mathrm{tot}}$ » est une
lecture incertaine : un $\mathfrak{T}$ (ou $\mathfrak{I}$) indicé
« total », pour l'ensemble des ordres totaux ; après
$\underline{\omega}(I\setminus\{i\})$, une formule biffée
($\{i\} \amalg \ldots$) n'est pas lue. La note marginale $\circledast$,
écrite de biais sur une dizaine de lignes courtes dans le coin supérieur
gauche, est d'une lecture très incertaine ; elle renvoie aux deux schémas du
bas de la page.}

$*$
\note{sous un trait tiré en travers du bas de la page, deux schémas de
correspondances, rattachés par l'astérisque à la note marginale.}

\[
  \begin{array}{ccc}
    V^* & \xrightarrow[\;\sim\;]{\;x \,\mapsto\, \mathrm{Cent}_G(x)\;} &
      \text{sous-groupes cycliques d'ordre 4 de } G = \mathfrak{S}_I \\[4pt]
    \text{l'unique élément} \neq 1 \text{ de } \Gamma \cap V &
      \longleftarrow\!\!\!\longrightarrow & \Gamma
  \end{array}
\]

\[
  \begin{array}{c}
    V^* \longleftrightarrow
      \text{ss-groupes cycliques d'ordre 4 de } G = \mathfrak{S}_I \\[4pt]
    x \longmapsto \mathrm{Norm}, \qquad
    D \longmapsto \text{élément non nul du centre de } D \\[4pt]
    \text{ss-groupes de Sylow de } G = \mathfrak{S}_I \ (D)
  \end{array}
\]
\note{les deux schémas sont redessinés : dans le second, trois flèches
relient en triangle $V^*$, les sous-groupes cycliques d'ordre 4 et les
$2$-sous-groupes de Sylow $D$ (flèche double entre ces deux derniers) ;
$x \mapsto \mathrm{Norm}$ va de $V^*$ vers les Sylow, et
$D \mapsto$ « élément non nul du centre de $D$ » en revient, sur un mot
surchargé illisible. Dans le premier, la page écrit
$x \mapsto \mathrm{Cent}_G(x)$ vers les sous-groupes cycliques
d'ordre 4 ; le centralisateur dans $\mathfrak{S}_4$ d'un élément de $V^*$
est d'ordre 8 (un $2$-Sylow), ce que le second schéma, qui passe par les
Sylow, semble corriger. Transcrit tel quel.}

\subsection{2. Structures de cube combinatoire.}

\page{7}
Le cube peut se définir, soit \uncertain{comme} le \uncertain{cas} \ill{}
des polygones \uncertain{réguliers}) par un système
\[
  (1) \qquad (S, A, R)
\]
\uncertain{sommets}, arêtes, incidence, satisfaisant les axiomes connus, soit
\uncertain{comme} utilisant l'\uncertain{antipodisme} $\underline{a}$ de (1)
en termes de
\[
  (2) \qquad
  \begin{array}{c}
    S \\
    \big\downarrow {\scriptstyle \text{degré } 2} \\
    D \\
    (\simeq S/\underline{a})
  \end{array}
  \quad \text{un des diagonales (de card.\ 2)}
\]
\marginal{(2) \ill{} \ill{} \ill{} \ill{} de card.\ $2$, \ill{} \ill{} $D$
\ill{} $V^*$ \ill{} diagonales)}
soit en une famille de deux ensembles (les \ill{} de card.\ $2$)
\uncertain{soit} \uncertain{dualement} par \ill{} de
\[
  (3) \qquad
  \begin{array}{c}
    A \\
    \big\downarrow {\scriptstyle \text{degré } 2} \\
    C \\
    (\simeq A/\underline{a})
  \end{array}
  \quad \text{un des codiagonales}
\]
\marginal{\uncertain{ou} ?}
comme une \struck{\ill{}} famille (\ill{} paires d'\struck{\ill{}} \ill{})
des codiagonales de card.\ $2$. Chacune des descriptions (1) (2) (3) est
valable pour décrire le cube --- là \struck{\ill{}} \uncertain{on} \struck{\ill{}}
\add{\uncertain{décrite} \uncertain{plus loin}}, les \struck{\uncertain{dimensions}}
\uncertain{entrant} \ill{} \ill{} l'espace, l'ens.\ \uncertain{possible}
\ill{} point de vue \ill{} « sommets » ; l'autre le point de vue
« arêtes », \uncertain{comme} \ill{} \ill{} \ill{} \uncertain{bien connus} dans
l'un et l'autre cas, c'est une \ill{} à quelle élément \ill{}
\uncertain{l'orientation} \uncertain{explicite}
\note{page d'une lecture très incertaine, surtout dans ses dernières lignes,
où la graphie est rapide et où une tache couvre le bas de la page. Les trois
présentations (1), (2), (3) et les deux flèches « degré 2 » sont sûres ; le
« $D$ » de (2) et le « $C$ » de (3) sont respectivement les diagonales et
les « codiagonales », quotients de $S$ et $A$ par l'antipodie
$\underline{a}$. La note marginale (2), écrite de biais dans la marge
gauche, n'est lue qu'à l'état de fragments.}

\page{8}\note{p.~4 de l'auteur.}
fixes qui porte la structure\ldots). On peut aussi décrire ces structures en
termes de $S$ et de $A$ par la donnée d'un groupe $R$ (« forme » de
$\mathbb{Z}/4\mathbb{Z}$) opérant sur $S$ resp.\ sur $A$, qui devient un
torseur sous $\Gamma$
\[
  \begin{aligned}
    &(2\ \text{bis}) \qquad S \text{ torseur sous } \Gamma \\
    &(3\ \text{bis}) \qquad A \text{ torseur sous } \Gamma
  \end{aligned}
\]
ce qui fait finalement \struck{\ill{}} \add{\uncertain{cinq}} descriptions
des \struck{\ill{}} structures de « cube ». Je vais expliciter les passages
de l'une de ces descriptions à l'autre.

\begin{itemize}
  \item[] $2\,\text{bis}) \Rightarrow 2)$, $3\,\text{bis}) \Rightarrow 3)$ :
    on prend l'unique élément $\underline{a}$ d'ordre 2 de $\Gamma$, qui est une
    \uncertain{involution} \uncertain{sans} pts fixes de $S$ (resp.\ $A$), et on
    pose
    \[
      D = S/\underline{a}, \quad \text{resp.} \quad C = A/\underline{a}.
    \]
  \item[] $2)$ \struck{bis} $\Rightarrow 1)$ : Soit
    $A \subset \mathfrak{P}_2(S)$ \add{$\subset \mathfrak{P}(S)$} défini par
    $A = \mathfrak{P}_2(S) \setminus \operatorname{Im} D$, c'est un ens.\ à
    4 éléments, \struck{qui} et on a une relation \uncertain{tautologique}
    $R \subset S \times A$.
  \item[] $3) \Rightarrow 1)$ : Soit $S \subset \mathfrak{P}_2(A)$ défini
    par $S = \mathfrak{P}_2(A) \setminus \operatorname{Im} C$, ens.\ à 4
    éléments et on a une relation tautologique $R' \subset A \times S$
    d'où ${}^{s}R' \overset{\mathrm{def}}{=} R \subset S \times A$.
  \item[] $1 \Rightarrow 2\,\text{bis}$, $1 \Rightarrow 3\,\text{bis}$ : On
    considère
    \[
      \mathfrak{D} = \operatorname{Aut} \underline{Q}, \qquad
      \underline{Q} = \{S, A, R\},
    \]
    d'ordre 8, \struck{$\mathfrak{D}_A\, V_Q$ 2-sous-gr.\ de Sylow de
    $\mathfrak{S}_S$ ou $\mathfrak{S}_A$} l'image inverse des
    $V_S \subset \mathfrak{S}_S$ et celle de $V_A \subset \mathfrak{S}_A$ sont
    deux sous-groupes \struck{\ill{}} de $\mathfrak{D}$, notés
    $\mathbb{F}_2^C \simeq V_S$ et $\mathbb{F}_2^D \simeq V_A$, \ill{}
    d'ordre 4, \struck{\ill{}} \ill{} \ill{} :
\end{itemize}

\[
  \begin{tikzcd}[column sep=small, row sep=small]
    & \mathfrak{S}_S \\
    \mathfrak{D} \arrow[ur] \arrow[dr] & \\
    & \mathfrak{S}_A
  \end{tikzcd}
  \qquad\qquad
  \begin{tikzcd}[column sep=small, row sep=small]
    \mathbb{F}_2^C \simeq V_S \arrow[r, hook] \arrow[d, hook] & \mathfrak{S}_S \\
    \mathfrak{D} \arrow[ur, no head] \arrow[dr, hook] & \\
    \mathbb{F}_2^D \simeq V_A \arrow[u, hook] \arrow[r, hook] & \mathfrak{S}_A
  \end{tikzcd}
\]
\note{les quatre dernières lignes de la page sont biffées d'autant de
traits ; on y lit, sans certitude, « $V_A$, \ldots\ $S$, $A$ \ldots\
$\mathbb{F}_2^D$ --- \ldots\ $V_Q$ \ldots\ il y a deux éléments d'ordre 4
\ldots\ \uncertain{engendrent} \ldots ». Le premier diagramme est à droite
dans le texte ; le symbole à sa source, surchargé plusieurs fois, se lit
$\mathfrak{D}$. Le second est dans la marge gauche, en face des lignes
$1 \Rightarrow 2\,\text{bis}$ ; le trait de $\mathfrak{D}$ vers
$\mathfrak{S}_S$ y porte un crochet d'inclusion mais pas de pointe.
« cinq » (descriptions) est ajouté au-dessus d'un mot biffé et n'est pas
sûr.}

\page{9}
\struck{\ill{} $\mathfrak{D}^+$ \ill{} sous-groupe cyclique d'ordre 4, et $S$
et $A$ sont \ill{} \ill{} \ill{} torseurs sous \ill{} $\mathfrak{D}^+$.}
\note{les trois premières lignes, dans un cadre ouvert à gauche, sont
barrées d'un grand trait en zigzag.}
groupe d'ordre 2 \struck{central} qui est le centre, et donc l'élément $\neq 1$
est l'\underline{antipodisme} $\underline{a}$.

Il y a dans $\mathfrak{D}$ exactement deux éléments $u, u'$ d'ordre 4, inverses
l'un de l'autre, engendrant l'unique sous-groupe cyclique d'ordre 4, noté
$\mathfrak{D}^+$ --- et $S$ et $A$ sont des torseurs sous ce même groupe.
L'antipodisme est l'unique élément d'ordre 2 de $\mathfrak{D}^+$.
\struck{Le quotient $\mathfrak{D}/\mathfrak{D}^+$ opère} \struck{$\mathfrak{D}^+
\simeq \mathfrak{D}/\{1, \underline{a}\}$}

Le quotient
$\mathfrak{D}/V_S$ \struck{\ill{}} $\simeq \mathbb{Z}/2\mathbb{Z}$ opère sur
$V_S$ par l'identité sur $\{1, \underline{a}\} \simeq \mathbb{F}_2 \subset V_S
\simeq \mathbb{F}_2^C$, et permute les deux points \struck{\ill{}} sur
$V_S \setminus \mathbb{F}_2 \simeq C$, et symétriquement
$\mathfrak{D}/V_A \simeq \mathbb{Z}/2\mathbb{Z}$ opère sur
$V_A \simeq \mathbb{F}_2^D$, \struck{\ill{}} l'opération par transposition sur
$D$. \struck{\ill{} \ill{}} \add{\ill{}}

Le quotient $\mathfrak{D}/\mathfrak{D}^+ \simeq \mathbb{Z}/2\mathbb{Z}$ opère
sur $\mathfrak{D}^+$ par l'automorphisme \struck{\ill{}} $x \mapsto x^{-1}$,
\struck{\ill{}} \add{en particulier}
$u \mapsto u' = u^{-1}$ i.e.\ l'opération induit la transposition sur
l'ens.\ $\simeq \underline{\omega}_Q$ des deux générateurs de $\mathfrak{D}^+$.
\marginal{$\mathfrak{D}^+ \simeq \mathbb{Z}/4\mathbb{Z} \wedge_{\{\pm 1\}}
\underline{\omega}_Q$}

On a \struck{\ill{}} \uncertain{finalement}
$\mathfrak{D}_{\mathrm{ab}} \simeq \mathfrak{D}/\{1, \underline{a}\}$, espace
vect.\ de rg~2 sur $\mathbb{F}_2$, \uncertain{avec} \uncertain{trois} \ill{}
$V_S$, $\mathfrak{D}^+$, $V_A$ $\longmapsto$ les images dans les trois
sous-groupes d'ordre 2 de $\mathfrak{D}^+_{\mathrm{ab}}$ \uncertain{sont}
\ill{} les trois seuls sous-groupes d'indice 2 de $\mathfrak{D}^+_{\mathrm{ab}}$.
On a une \uncertain{bonne} \ill{} de $\mathfrak{D}^+_{\mathrm{ab}}$
\[
  \mathfrak{D}^+_{\mathrm{ab}} \simeq \mathbb{F}_2 \times \mathbb{F}_2
\]
définie par la condition que
\marginal{$\mathfrak{D}$ opère sur $C$ \uncertain{comme} $D$ par
\ill{} \ill{} \ill{}, \ill{} $C \amalg D$ \ill{} \ill{} \ill{} $=3$, $D$
opère sur $C$, $D$ opère sur $C$ \ill{} i.e.\ $c \in C$, $d \in D$
$\Rightarrow [c, d] = \underline{a}$) d'où \ill{} $\mathfrak{D} \to
\mathfrak{S}_C \times \mathfrak{S}_D$ \ill{} $\simeq \mathbb{F}_2 \times
\mathbb{F}_2$ \ill{} \ill{} \ill{} \ill{} $\mathfrak{D}/\{1, \underline{a}\}
\simeq V_S/\mathbb{F}_2 \times V_A/\mathbb{F}_2$
\underline{NB} $\mathfrak{D}_{\mathrm{ab}} \simeq
\mathbb{F}_2^C \times \mathbb{F}_2^D$}
\note{la note marginale, écrite en travers du coin inférieur gauche sur une
dizaine de lignes obliques et en partie surchargée de grands traits, est
d'une lecture très incertaine ; ses formules sont plus sûres que ses mots.
Le texte parle tantôt de $\mathfrak{D}^+_{\mathrm{ab}}$, tantôt de
$\mathfrak{D}_{\mathrm{ab}}$ ; le « $+$ » en exposant est net aux deux
dernières occurrences, bien que le sens (quotient de $\mathfrak{D}$ par son
centre, de rang 2) soit celui de $\mathfrak{D}_{\mathrm{ab}}$. Transcrit tel
quel.}

\page{10}\note{p.~5 de l'auteur.}
\[
  \left\{
  \begin{array}{lll}
    \text{image inverse de} & \mathbb{F}_2 \times \{0\} & \leadsto V_S \\
    \text{---\,---} & \{0\} \times \mathbb{F}_2 & \leadsto V_A \\
    \text{---\,---} & \text{diagonale de } \mathbb{F}_2 \times \mathbb{F}_2
      & \leadsto \mathfrak{D}^+
  \end{array}
  \right.
\]
on en \uncertain{conclut}
\[
  \left\{
  \begin{array}{lll}
    \text{image inverse de} & (1, 0) & \leadsto C \subset \mathfrak{D}^* \\
    \text{---\,---} & (0, 1) & \leadsto D \subset \mathfrak{D}^* \\
    \text{---\,---} & (1, 1) & \leadsto \underline{\omega}_Q \subset \mathfrak{D}^*
  \end{array}
  \right.
\]
\note{les $V_S$, $V_A$ du premier tableau sont écrits sur un
$\mathfrak{D}$ surchargé ; les tirets sont les traits de répétition de
l'auteur (« idem »). Dans le second tableau, $\mathfrak{D}^*$ est lu pour
un $\mathfrak{D}$ muni d'un exposant peu net, sans doute $\mathfrak{D}^*$
au sens des éléments $\neq 1, \underline{a}$.}

Les sous-groupes d'ordre deux de $\mathfrak{D}$ correspondent aux éléments
d'ordre 2, il y a \uncertain{cinq} :

les éléments
\begin{itemize}
  \item[a)] l'antipodisme $\underline{a}$, engendrant le centre
    $\{1, \underline{a}\}$
  \item[b)] Les deux éléments de $C = V_S \setminus \{1, \underline{a}\}$,
    engendrant resp.\ les deux droites-coordonnées de
    $V_S \simeq \mathbb{F}_2^C$
  \item[c)] Les deux éléments de $D = V_A \setminus \{1, \underline{a}\}$,
    engendrant resp.\ les deux droites coordonnées de
    $V_A \simeq \mathbb{F}_2^D$.
\end{itemize}
(NB les \ill{} \uncertain{autres} 3 éléments de $\mathfrak{D}^*$ sont
d'ordre 4 :
\[
  \mathfrak{D} \simeq
  \underbrace{\{1\} \amalg \{a\} \amalg C \amalg D}_{\text{dic.\ de }
  \mathfrak{D} \text{ en classes de conjugaison}} \amalg \underline{\omega}_Q
  \qquad )
\]
\note{sous la formule, deux accolades se chevauchent ; la seconde porte
« \ldots\ orbites sous $\mathrm{Int}(\mathfrak{D})$ » ; leur étendue exacte
n'est pas sûre.}
On a donc fait la liste complète des sous-groupes de $\mathfrak{D}$ différents
de $\{1\}$ et de $\mathfrak{D}$ lui-même, il y en a trois d'ordre 4 i.e.\
d'indice 2, savoir $V_S$, $V_A$ et $\mathfrak{D}^+$, et cinq d'ordre 2 sauf le
centre $\{1, \underline{a}\}$, et les \add{2 paquets de 2}
\struck{\uncertain{4}} sous-groupes d'ordre 2 distincts de ce dernier,
contenus respectivement dans $V_S$ et dans $V_A$.
\marginal{NB. On \uncertain{retrouve} par \uncertain{ceci} les éléments
\uncertain{réunis} \ill{} \ill{} d'un des \ill{} \ill{} \ill{} leurs
\ill{} \ill{} \ill{} \ill{} \ill{} $\mathfrak{D}$ \ill{} \ill{} $V_A$ et
$V_S$ \ill{} \ill{} $\mathfrak{D}$, distingués et \uncertain{quotients}
$\neq \{1\}, \mathfrak{D}$ \ill{} $\{1, \underline{a}\}$) \ill{}
$V_S$, $V_A$, $\mathfrak{D}^+$, \ill{} \ill{} \ill{}}

On a associé à un cube $Q = (S, A, R)$ trois ensembles \uncertain{bien-gardés}
de cardinal 2, savoir
$\{$ \uncertain{$D$}, \uncertain{$C$}, $\underline{\omega}_Q$
\note{la dernière ligne de la page s'arrête sur cette énumération, que
la page~11 poursuit (« en plus des deux diagonales et des deux
codiagonales\ldots »). Les deux premiers symboles ont l'air d'un
$\mathfrak{B}$ et d'un $S$ ; la suite du texte ($\underline{\omega}_S \simeq
C$, $\underline{\omega}_A \simeq D$) désigne $D$ et $C$. Le mot qualifiant les
trois ensembles n'est pas sûr. La note marginale, oblique dans la marge gauche, n'est lue que par
fragments.}

\page{11}
(en plus des deux diagonales et les deux codiagonales elles-mêmes, qui font
aussi \uncertain{partie} \uncertain{des} ensemble : deux éléments, mais
groupés par paires « indistinguables » \ldots). On a des bijections
canoniques
\[
  \underline{\omega}_S \simeq C, \qquad \underline{\omega}_A \simeq D
\]
provenant du fait que\ldots
\[
  \underline{\omega}_S \simeq \underline{\omega}_{V_S^*}
  = \underline{\omega}_{\{\underline{a}\} \amalg C}
  \simeq \underline{\omega}_C \simeq C
\]
la convention pour orienter un ensemble : trois éléments (\uncertain{ici}
$V_S^*$) dont un élément $x$ (ici $\underline{a}$) est fixé, en termes de
\uncertain{choix} d'un élément $y$ des \uncertain{deux} \uncertain{autres}
[ici $V_S^* \setminus \{a\} = C$] étant de prendre l'orientation
$(a, x, y)$ \struck{($\ldots$ $y \neq a, x$ $\ldots$ $= V_S^* \setminus
\{a, \ldots$)} i.e.\ $\{y\} = J \setminus \{a, x\}$.
\note{les lettres $a$, $x$, $y$ de ce passage ne s'accordent pas tout à
fait : « un élément (ici $\underline{a}$) est fixé » et un « $x$ » écrit
au-dessus ; transcrit tel quel. Le $J$ de la dernière formule est de
l'auteur.}

On a une relation remarquable
\[
  (*) \qquad \underline{\omega}_Q \wedge C \wedge D \simeq \mathbb{1}_{\mathbb{F}_2}
\]
\marginal{i.e.\ une restriction du groupe structural $\mathbb{F}_2^3$ de
$\underline{\omega}_S \times C \times D$ au sous-groupe $V_3(\mathbb{F}_2)$
des $x, y, z$ tels que $x + y + z = 0$}
qui provient du fait que $\underline{\omega}_S$, $C$, $D$ sont les images
inverses, dans l'extension $\mathfrak{D}$ de
$\mathfrak{D}/\{1, \underline{a}\} = \mathfrak{D}_{\mathrm{ab}}$ par
$\{1, \underline{a}\}$, des trois éléments non nuls de
$\mathfrak{D}_{\mathrm{ab}}$, dont la \struck{produit} somme est nulle.
\add{Pour} des éléments $\varpi$, $c$, $d$ de
$\underline{\omega}_Q \times C \times D$ \struck{\ill{}} \ldots
\[
  \varpi \wedge c \wedge d = 0 \iff \varpi c d = 1 \text{ dans } \mathfrak{D}
\]
[Attention, cette relation \struck{serait} \uncertain{invariante} si \ill{}
faisait une transposition sur l'ordre des facteurs --- il y a donc
\struck{\ill{}} des \struck{\ill{}} le choix de l'iso.\ \uncertain{ces} une
question de convention.
\note{les deux « $\underline{\omega}_S$ » de ce passage (dans la note
marginale et dans « qui provient du fait que ») sont écrits ainsi, là où la
formule $(*)$ porte $\underline{\omega}_Q$.}

Si on passe du cube $Q$ au cube dual $Q'$, le choix de l'iso
\[
  \underline{\omega}_{Q'} \wedge C' \wedge D' \simeq \mathbb{1}
  \ (\simeq \underline{\omega}_Q \wedge D \wedge C) \simeq \mathbb{1}
\]
est \underline{opposé} de celui pour $Q$\ldots
\note{dans la dernière formule, un mot biffé devant $C'$ et un autre
devant le premier $\mathbb{1}$.}

\page{12}\note{p.~6 de l'auteur.}
\underline{NB} L'ensemble produit : 4 éléments
\[
  C \times D
\]
peut s'identifier : l'ensemble quotient $R/\underline{a}$ des repères du cube
$Q$, modulo antipodisme
\[
  C \times D \simeq R/\underline{a} \qquad
  R = \operatorname{Rep}(Q) \simeq \operatorname{Rep}(Q^\circ)
\]
tandis que $\underline{\omega}_Q$ s'identifie à \struck{l'ensemble} un
ensemble de deux parties de $R$ (« $R^+$ » et « $R^-$ ») complémentaires,
l'une de l'autre, et stables chacune par l'antipodisme $\underline{a}$ (qui
conserve l'orientation), ce qui fait que le choix d'un élément de
$R/\underline{a}$ définit canoniquement une orientation de $Q$,
\struck{\ill{}} i.e.\ un des deux \struck{\ill{}} ensembles de repères
$R'$, $R''$. Pour vérifier \struck{que} si cette description
\[
  C \wedge D \simeq \underline{\omega}_Q \quad \text{i.e.\ d'une relation}
  \quad \varpi c d = 1
\]
d'une iso. $\underline{\omega}_Q \wedge C \wedge D \simeq \mathbb{1}$
correspond : \ill{} \struck{$\ldots \wedge \ldots$}, il faut expliciter de
bout en bout la convention qui nous \struck{\ill{}} conduit aux plongements
\[
  C \hookrightarrow V_S \setminus \{1, \underline{a}\} \subset \mathfrak{D},
  \qquad D \hookrightarrow V_A \setminus \{1, \underline{a}\} \subset \mathfrak{D},
  \qquad \underline{\omega}_Q \subset \mathfrak{D}^+ \subset \mathfrak{D}
\]
Si $c$ est une \add{codiagonale i.e.\ une} paire de côtés opposés du cube,
je vais lui associer l'automorphisme $\sigma$ du cube qui fixe ces deux
côtés et qui échange les deux autres i.e.\ la symétrie p.r.\ à l'axe
transverse. De même, pour une diagonale, je lui associe la symétrie par
rapport à cette diagonale. Pour une orientation définie par un repère
\ill{}

\note{dans la marge gauche, trois petites figures : en haut, un carré de
sommets $s$, $s'$, d'arêtes $a$ (en bas) et $a'$ (en haut) épaissies, et de
diagonale $d$ tracée de $s$ à $s'$ --- un repère du cube, vu sur une face ;
en bas, un carré dont deux côtés opposés, épaissis, portent la marque $c$,
avec en pointillé l'axe transverse qui les coupe en leur milieu ; à droite,
un carré en pointillé dont la diagonale, prolongée en pointillé au-delà des
sommets, porte la marque $d$ à ses deux extrémités.}

\page{13}
\ldots\ je lui associe la rotation du cube qui transforme $s$ en la deuxième
extrémité de $a$. On vérifie que avec ces conventions, la formule
$\varpi c d = 1$ est bien la formule qui exprime que $R_\varpi \subset R$
contient l'image inverse par $R \to R/\underline{a}$ de l'élément $(c, d)$ de
$R/\underline{a}$. Le fait que si on remplace la $\varpi$-relation par le cube
dual, on semble obtenir la relation opposée, provient du fait que si on
identifie le $\varpi$-ensemble $R$ à l'ens.\ des repères de $Q$ \underline{et}
de $Q^\circ$, \struck{\ill{}} \add{et} le groupe $\mathfrak{D}$ à
$\operatorname{Aut} Q$ \underline{et} à $\operatorname{Aut}(Q^\circ)$ (ce qui
était \uncertain{possible} \uncertain{naturellement} conséquence que d'une
seule façon, i.e.\ ce n'est pas une question de choix de conventions\ldots)
alors une $\varpi$-\struck{partie} \add{classe} $Q'$ dans $R$
\struck{\ill{}} (pour la relation d'équivalence de l'orientation),
considérée comme définissant une orientation de $Q$ \underline{et} une
orientation de $Q^\circ$, définit deux orientations \underline{opposées},
en ce sens que les orientations de rotations qu'il définit sur un repère
$(s, a)$, suivant qu'il est regardé comme un rep.\ de $Q$ ou rep.\ de
$Q^\circ$, sont \underline{opposées} l'une de l'autre : transformer $s$ en
$s'$ ou $a$ en $a'$ se font par deux rotations \underline{opposées}.

\note{dans la marge gauche, deux petites figures. En haut, un carré en
pointillé dont un côté, plein, est marqué $a$, avec le sommet $s$ à son
extrémité inférieure et une flèche courbe tournant autour de l'autre
extrémité : la rotation qui envoie $s$ sur l'autre extrémité de $a$. En bas,
un carré de sommets $s'$ et $s$ en haut, l'arête supérieure $a$ (de $s'$ à
$s$) épaissie, l'arête droite marquée $a'$, et une flèche courbe au-dessus
de $s$ : la rotation qui transforme $s$ en $s'$ et $a$ en $a'$.}

\page{14}\note{p.~7 de l'auteur.}
Considérons le groupe $\mathfrak{D}$ en tant que groupe « abstrait », il est
canoniquement extension de $\mathfrak{D}_{\mathrm{ab}} = W$
($\simeq \mathbb{F}_2 \times \mathbb{F}_2$ non canoniquement) par
$\mathbb{F}_2$. Dans $\mathfrak{D}_{\mathrm{ab}}$ il y a un élément non nul
distingué $\varpi$, \add{dont l'image} \struck{\ill{} \ill{}} inverse est
l'ens.\ $\underline{\omega}_Q$ des deux éléments d'ordre 4 de $\mathfrak{D}$.
On voit aisément que les deux autres éléments non nuls de
$\mathfrak{D}_{\mathrm{ab}}$ sont interchangés par le groupe des autom.\ de
$\mathfrak{D}$, on trouve donc
\[
  1 \to \underset{\substack{\text{autom.\ de l'ext.\ } \mathfrak{D} \text{ de}\\
  W \text{ par } \mathbb{F}_2 \text{ qui}\\ \text{induisent l'identité}\\
  \text{sur } W \text{ et } \mathbb{F}_2}}{\check{W}}
  \to \operatorname{Aut} \mathfrak{D} \to
  \underset{\substack{\text{interchange}\\ \text{les deux}\\ \text{éléments de}\\
  W \setminus \{0, \varpi\}}}{\underbrace{\pm 1}} \to 1
\]
\marginal{\underline{NB} on a un iso.\ can.\ $\check{W} \simeq W$ (valable
pour \uncertain{tt vectoriel} de rg~2 sur $\mathbb{F}_2$)}
en fait $\operatorname{Aut} \mathfrak{D}$ est produit $\frac12$ direct de
$\pm 1$ par $\check{W}$, $\pm 1$ opérant par échange des \add{2} éléments non
\struck{\ill{}} privilégiés de $\check{W} \simeq W$. Quand on tient compte que
$\mathfrak{D}$ provient d'un cube, donc il y a choix d'un élément
$\underline{c}$ de $W \setminus \{0, \varpi\}$ (correspondant
\struck{\ill{}} à $C = V_S \setminus \{1, \underline{a}\}$), l'autre élément
étant $\underline{d}$ (correspondant à
$D = V_A \setminus \{1, \underline{a}\}$). Les autom.\ de $\mathfrak{D}$ qui
respectent cette structure forment le sous-groupe
$\check{W} \simeq W$ de $\operatorname{Aut} \mathfrak{D}$, formé des autom.\
qui induisent l'identité sur le quotient $W$. Ce sont aussi les autom.\
intérieurs, qui forment bien le groupe
$W \simeq \mathfrak{D}/\underbrace{\{1, \underline{a}\}}_{\text{centre}}$.
(\add{l'homom.} $\mathfrak{D} \xrightarrow{\;\mathrm{int}\;}
\operatorname{Aut} \mathfrak{D}$ n'est pas \struck{homom} injectif), on voit
que la donnée d'un groupe $\mathfrak{D} \simeq \operatorname{Aut} Q$ (qui
revient : la donnée de $C$ et de $D$, et plus \ill{}
\struck{\ill{} $\mathfrak{D} = \mathbb{F}_2^{C} \times_{\mathbb{F}_2}
\mathbb{F}_2^{D}$})
\note{la phrase de la fin de page est très abrégée. La parenthèse
« (qui revient : \ldots) » se referme à droite sur la dernière ligne, biffée
d'un double trait, où on lit $\mathfrak{D} = \mathbb{F}_2^{C}
\times_{\mathbb{F}_2} \mathbb{F}_2^{D}$ ($\mathbb{F}_2$ écrit sous le signe
$\times$) ; au-dessus, « et plus » est suivi d'un symbole muni d'un exposant
$(\underline{a}, s)$, non lu. Au début de la page, le mot qui qualifie
l'image inverse de $\varpi$ est inséré au-dessus d'un passage biffé.}

\page{15}
ne peut suffire à reconstruire $Q$ à isom.\ unique près --- mais seulement
à $Q$ « modulo antipodisme ». Considérons le cube épinglé $Q_0$, et le groupe
de ses automorphismes
\[
  \mathbb{D}_4 = \operatorname{Aut}(Q_0)
\]
on a la suite exacte canonique
\[
  1 \to \mathbb{F}_2 \to \mathbb{D}_4 \to \mathbb{F}_2^{2} \to 1
\]
et la donnée d'un tel $\mathfrak{D}$ revient à celle d'un torseur sous
$\mathbb{F}_2^{2}$ \add{le quotient de $\mathbb{D}_4$} (i.e.\ de deux ens.\ à
deux éléments (ce sont en fait $D$ et $C$, dans cet ordre\ldots)), tandis que
celle d'un cube revient à la donnée d'un torseur sous $\mathbb{D}_4$
lui-même.

On a une opération \add{(par autom.\ int.)} de $\mathfrak{D}^+$ sur $D$
\add{$\hookrightarrow \mathfrak{D}^+$}, la donnée de $Q$ (correspondant au
groupe d'autom.\ $\mathfrak{D}$) \struck{\ill{}} revient à celle d'un
relèvement de \struck{\ill{}} l'opération \struck{\ill{}} en \struck{les} une
opération \add{(d'un ens.\ à 4 élts sous $\mathfrak{D}$)} transitive sur un
ens.\ $S$, avec $S/\underline{a} \simeq D$. Description duale via
$C \hookrightarrow \mathfrak{D}^+$\ldots
\note{« $\hookrightarrow \mathfrak{D}^+$ » est écrit au-dessus de $D$, et le
« $+$ » en exposant est net, bien que $D \subset V_A$ ; de même dans
« $C \hookrightarrow \mathfrak{D}^+$ ». Transcrit tel quel.}

\underline{Remarque} : Le groupe $\operatorname{Aut}(\mathbb{D}_4)$ est isom.\
à $\mathbb{D}_4$ (mais pas via autom.\ intérieurs, bien sûr) ce qui indique
qu'il y a \add{peut-être} une équivalence de groupoïdes
\[
  (\text{Cubes}) \xrightarrow{\;\approx\;} (\text{groupes isom.\ à }
  \mathbb{D}_4)
\]
(mais qui n'est pas donnée par \struck{\ill{}} $Q \mapsto \operatorname{Aut}(Q)$ !),
qui mériterait d'être comprise\ldots
\note{$\mathbb{D}_4$ est écrit, à sa première occurrence, après un
$\mathfrak{D}_4$ biffé ; c'est la notation de la page pour le groupe du cube
épinglé.}

\page{16}\note{p.~8 de l'auteur.}
Si $\mathfrak{D}$ est un groupe isomorphe à $\mathbb{D}_4$, l'ens.\
$\mathfrak{S}$ des éléments d'ordre 2 \add{$\neq \underline{a}$ (où
$\underline{a}$ =} \struck{\ill{}} l'unique élément central $\neq 1$) est à
quatre éléments, \struck{\ill{} \ill{} déjà classés par translation}
\add{\ill{} une involution} sans points fixes par
\[
  x \longmapsto x \cdot \underline{a} \overset{\mathrm{def}}{=} x'
\]
donc à $\mathfrak{D}$ correspond un cube $Q_{\mathfrak{D}}$, dont l'ens.\ des
sommets est $\mathfrak{S}$, et l'\struck{\ill{} des} \add{\ill{}}
l'antipodisme sur $S$ est l'involution précédente. Mais le foncteur de
groupoïdes $\mathfrak{D} \mapsto Q_{\mathfrak{D}}$
\[
  (\text{groupes isom.\ à } \mathbb{D}_4) \longrightarrow
  (\text{cubes combinatoires})
\]
\ill{} \ill{} une équivalence de groupoïdes, on \ill{} reconstitue
\add{$\mathfrak{D}$} à partir de $(\mathfrak{S}, \underline{a}_{\mathfrak{S}}
\colon x \mapsto x')$ \struck{\ill{}} par la formule
\[
  \mathfrak{D} = \Bigl\{ \underbrace{x \in \mathfrak{S},\ \underline{a}}_{
  \text{générateurs}} \Bigm| \underbrace{[x, y]}_{\substack{\text{commutateur}\\
  xyx^{-1}y^{-1}}} = \underline{a} \text{ si } y \neq x, x' \Bigr\}
\]
\marginal{\underline{NB} Si \uncertain{$Q$} est un cube, \struck{\ill{}} le
groupe $\mathfrak{D} = \operatorname{Aut}(Q)$, le carré associé à ce groupe, est
le carré dont l'ens.\ des sommets est $C \amalg D$, l'antipodisme dessus
\struck{\ill{}} \ill{} $C$ et $D$, i.e.\ les $C$ et $D$. On trouve donc
\ill{} \ill{} \struck{\ill{}} distinguées\ldots}
\note{la note marginale, écrite en oblique dans la marge gauche et reliée
par un trait pointillé à l'accolade de la formule, n'est lue que par
fragments ; on y lit « carré », non « cube ». L'ensemble des éléments
d'ordre 2 est noté par une lettre ronde, lue $\mathfrak{S}$, et le $S$ de
« l'antipodisme sur $S$ » est droit. Le début de la phrase « \ldots\ une
équivalence de groupoïdes » est surchargé d'ajouts entre les lignes et
n'est pas lu : on ne peut dire s'il affirme ou nie l'équivalence.}

Si on part, pour décrire la structure de cube, de l'ens.\ des sommets $S$,
cette structure est définie par la donnée d'un ens.\ \add{$S$} à 4 élément
--- i.e.\ d'un torseur sous un vectoriel $V_S$ de rang 2 sur $\mathbb{F}_2$
--- et de une \struck{involution} sans pts fixes de $S$, i.e.\ d'un
$\underline{a}_S \in V_S^{*}$ \struck{Mais \ill{}} \add{(i.e.\ une}
\struck{$V_S = V_S^{*}$} partition de type $(2,2)$ de $S$). Mais pour

\page{17}
un vectoriel \add{$V$} de rg~2 sur $\mathbb{F}_2$, la donnée d'un
$\underline{a} \in V^{*}$ revient à la donnée d'une partie de $V^{*}$ qui soit
une base (à savoir $V^{*} \setminus \{\underline{a}\}$) ; donc la donnée d'un
\add{couple $(V, \underline{a})$ d'un} vectoriel $V$ de rg~2 sur
$\mathbb{F}_2$, et d'un élément $\underline{a}$ de $V^{*}$, revient à la donnée
d'un ens.\ $C$ à deux éléments, par $V \simeq \mathbb{F}_2^{C}$,
$\underline{a}$ l'élément diagonal de $\mathbb{F}_2^{C}$. Ainsi, la donnée
d'un cube $Q$ \struck{\ill{} \ill{}} (via l'ens.\ des sommets) équivaut à
celle d'un couple $(C, S)$, où $C$ \struck{\ill{}} est un ens.\ à deux
éléments, et $S$ est un torseur sous $\mathbb{F}_2^{C}$. Mais la donnée d'un
tel torseur revient à celle d'une famille de \add{$\mathbb{F}_2$-}torseurs
indexée par $C$, i.e.\ d'un revêtement de degré 2
\[
  A \longrightarrow C,
\]
\add{(\ill{} de $C$)} i.e.\ d'un ens.\ $A$ à quatre éléments, et muni d'une
involution sans pts fixes (qui donne le quotient vers $C$), par la formule
\[
  S \simeq \prod_{c \in C} A_c = \text{Ens des sections de } A
  \text{ sur } C
\]
qui donne en même temps un plongement
\[
  S \hookrightarrow \mathfrak{P}_2(A)
\]
d'où la relation d'incidence entre $S$ et $A$\ldots
\marginal{ici les \struck{deux} --- $A_c$ ($c \in C$) \ill{}
\uncertain{comme} \ill{} \uncertain{quotients} de $S$, par les relations qui
distinguent l'élément \add{$c \in C$} \struck{\ill{}} de $V^{*}$
\uncertain{considéré} comme involution ss pts fixes de $S$\ldots\ Mais je
préfère ici prendre $A$ \struck{\ill{}} \ill{} \ill{}}
\note{la note marginale, oblique sur une dizaine de courtes lignes dans le
coin inférieur gauche, n'est lue que par fragments. Au début de la page,
l'$\underline{a}$ de « $\underline{a} \in V^{*}$ » est écrit devant un
$\mathbb{F}_2^{*}$ surchargé en $V^{*}$. Dans la formule
$S \simeq \prod_{c \in C} A_c$, un premier symbole biffé précède $A_c$.}

\page{18}\note{p.~9 de l'auteur.}
Dualement, \struck{$A$ \ill{}} la donnée d'un cube-arêtes
$A$ s'interprète comme la donnée d'un torseur sous un groupe
$\mathbb{F}_2^{D}$, ce qui équivaut à la donnée d'une famille de
\add{$\mathbb{F}_2$-}torseurs $(S_d)_{d \in D}$, \struck{i.e.\ \ill{} \ill{}
ens.\ $S$}
\[
  A \simeq \prod_{d \in D} S_d \qquad (\simeq \text{sections de } S \text{ sur } D)
\]
--- la donnée de cette famille équivaut à la donnée d'un ensemble (à quatre
éléments) $S$ et d'un morphisme de degré 2
\[
  S \longrightarrow D
\]
qui \ill{} la formule précédente donne
\[
  A \hookrightarrow \mathfrak{P}_2(S)
\]
et la relation d'incidence entre $S$ et $A$\ldots
\note{le « Dualement » initial est écrit sur un premier mot surchargé ; le
début de phrase biffé qui suit n'est pas lu.}

\underline{Formulaire}. Rappelons, pour un cube $Q$, les trois ensembles à
deux éléments
\[
  \underset{\text{codiagonales}}{C}, \quad
  \underset{\text{diagonales}}{D}, \quad
  \underset{\text{orientations}}{\underline{\omega}}
\]
plus les deux paquets de deux ens.\ à deux éléments (dans \ill{} autres ens.\
à deux éléments)
\[
  S \xrightarrow{\;\deg.\,2\;} D, \qquad A \xrightarrow{\;\deg.\,2\;} C
\]
\[
  (D \simeq S/\underline{a}_S). \qquad (C \simeq A/\underline{a}_A)
\]
On a le formulaire d'iso.\ can.
\note{le symbole du troisième ensemble, souligné, n'est pas net et sa lecture $\underline{\omega}$ est incertaine ; la
légende « orientations » placée dessous désigne $\underline{\omega}_Q$.}

\page{19}
\[
  \begin{array}{ll}
    (1) & \underline{\omega}_S \simeq C, \quad \underline{\omega}_A \simeq D \\[6pt]
    (2) & \underline{\omega} \wedge C \wedge D \simeq \mathbb{1}_{\mathbb{F}_2}
      \quad \text{i.e.} \quad
      \left\{ \begin{array}{l}
        \underline{\omega} \simeq C \wedge D \\
        C \simeq D \wedge \underline{\omega} \\
        D \simeq \underline{\omega} \wedge C
      \end{array} \right. \\[18pt]
    (3) & D \simeq S/\underline{a}_S, \quad C \simeq A/\underline{a}_A \\[6pt]
    (4) & S \simeq \prod_{c \in C} A_c, \quad A \simeq \prod_{d \in D} S_d \\[6pt]
    (5) & D \simeq \bigwedge_{c \in C} A_c, \quad C \simeq \bigwedge_{d \in D} S_d
  \end{array}
\]
ces \add{deux} dernières relations proviennent de
\[
  \underline{\omega} \simeq \Bigl(\bigwedge_{c \in C} A_c\Bigr) \wedge C
  \quad \text{d'où} \quad
  \underbrace{\underline{\omega} \wedge C}_{D \text{ par } (2)} \simeq
  \bigwedge_{c \in C} A_c .
\]
\note{le formulaire est encadré. Dans (3), le $D$ est écrit sur un $C$.}

\subsection{3. Compléments sur les 2-sous-groupes d'un groupe $G = \mathfrak{S}_I$ (card $I = 4$).}

\note{dans le titre, « 2-s-groupes » est écrit sur une première lettre, et
« de Sylow » qui le suivait est biffé. En marge, en oblique : « À
\uncertain{compléter} par les 3-sous-groupes de Sylow ».}
Soit $V$ le sous-groupe distingué $\simeq \mathbb{F}_2^{2}$, de sorte que
$V^{*} \simeq \mathfrak{P}_{2,2}(I)$. On a

\[
\begin{tikzcd}[column sep=large, row sep=huge]
  V^{*} \arrow[rr, leftrightarrow, "\sim"] \arrow[dr, leftrightarrow, "2"'] & &
  \text{ss-gr.\ cycliques } \mathfrak{D}^+ \text{ d'ordre 4 de } G \arrow[dl, leftrightarrow, "\sim"] \\
  & \text{2-ss-gr.\ de Sylow de } G &
\end{tikzcd}
\]

\begin{itemize}
  \item[] à gauche de $V^{*}$ : « Ens.\ des structures de cubes $Q$
    $\simeq V^{*}$ (si $Q$ correspond à $x$, \ill{} \ill{} $\mathfrak{D}$,
    \ill{} $\operatorname{Aut} Q = \mathfrak{D}$) » ;
  \item[] sous la flèche horizontale : « pour tt ss gr.\ cyclique
    $\mathfrak{D}^+$ il y a un unique élt d'ordre 2, qui est dans $V^{*}$, et
    $\mathfrak{D}^+$ en est le centralisateur » ;
  \item[] le long de la flèche de gauche : « pour tt sous-groupe de Sylow
    $\mathfrak{D}$, il y a un unique élément $x$ central $\neq 1$, qui est
    celui qui est dans $V^{*}$, et $\mathfrak{D}$ est le normalisateur ou le
    centralisateur de $x$ » ;
  \item[] le long de la flèche de droite : « tout sous-groupe de Sylow
    $\mathfrak{D}$ a un unique sous-groupe cyclique d'ordre 4,
    $\mathfrak{D}^+$, et il en est le normalisateur ».
\end{itemize}
\note{le diagramme est redessiné : un triangle à flèches doubles, dont les
sommets et les légendes sont écrits en toutes lettres ; les légendes sont
rendues ci-dessus dans l'ordre des flèches. L'étiquette de la flèche de
gauche est lue « 2 » ; le chiffre initial du sommet du bas ressemble à un
« 3 », mais les groupes d'ordre 8 visés sont les 2-sous-groupes de Sylow. La
légende de gauche est séparée de celle de la flèche par un trait.}

\page{20}\note{p.~10 de l'auteur.}
Quels sont les sous-groupes de $G = \mathfrak{S}_I$ \struck{\ill{}},
$\neq e$ et $\mathfrak{S}_I$ ?

A) Sous-groupes d'ordre une puissance de 2, i.e.\ contenus dans un
2-sous-groupe de Sylow $\mathfrak{D}$.

\struck{Sous-groupes cycliques correspondant aux}

1) Sous-groupes d'ordre 2, correspondant aux éléments d'ordre 2 :
\begin{itemize}
  \item[a)] les \add{trois} éléments de $V^{*}$, qui forment une orbite
    \struck{sous} $\operatorname{Int} G$
  \item[b)] les six transpositions ou \uncertain{réflexions}, qui forment une
    autre orbite sous $\operatorname{Int} G$, et qui se groupent en trois
    paquets (indexés par $V^{*}$) de deux, par la relation
    $\tau\tau' \in V^{*}$, ou par celle d'être contenus dans un même
    \add{2-}sous-groupe de Sylow. \struck{cycliq}
\end{itemize}
\marginal{i.e.\ $\tau$, $\tau'$ \ill{} \ill{} $\mathfrak{S}_I$ \ill{}}

2) Sous-groupes d'ordre 4
\begin{itemize}
  \item[c)] \add{les trois} \struck{Sous}-groupes cycliques d'ordre 4, qui
    \struck{correspondent \ill{}} \add{1-1 aux trois} sous-groupes de Sylow
    i.e.\ les trois structures de cube sur $I$.
  \item[d)] Sous-groupes non cycliques d'ordre 4, donc isom.\ à
    $\mathbb{F}_2 \times \mathbb{F}_2$ --- il y a d'abord $V^{*}$
  \item[e)] puis \struck{deux} ceux qui ne sont pas distingués, chacun contenu
    dans un unique 2-Sylow \struck{\ill{}} \add{un} \ill{} groupe de Sylow,
    dans lequel il y en a un \add{(en plus de $V$, qui y est contenu)} --- il
    y en a en tout \struck{six} \add{trois} sous-groupes
    \struck{$\mathbb{F}_2 \times \mathbb{F}_2$} $(2,2)$ non distingués, qui
    forment une classe de conjugaison \struck{distincte, formant deux \ill{}
    de trois classes \ill{} savoir les $V_S \subset \mathfrak{D}$ et les
    $V_A$}, correspondant \ill{}
\end{itemize}

3) Sous-groupes d'ordre 8 --- ce sont les
\begin{itemize}
  \item[f)] trois sous-groupes de Sylow de $G$
\end{itemize}
\note{dans la marge gauche, en face des items, deux colonnes : les
nombres d'éléments $3$, $6$, $3$, $1$, $3$, $3$ (a à f), et les ensembles qui
les indexent : $V^{*}$, $\widetilde{V}^{*}$, $V^{*}$, ---, $V^{*}$, $V^{*}$.
En d), la page écrit $V^{*}$, là où le compte « 1 » et le sens désignent le
sous-groupe $V$ lui-même ; transcrit tel quel. En e), les ajouts entre les
lignes et les passages biffés se chevauchent ; l'ordre de lecture retenu
n'est pas sûr.}

\end{document}
